\documentclass{article}
\usepackage{amsmath,amssymb,url}
\title{Settlement, retained costs and selective recovery at a compound cut}
\author{}\date{}
\begin{document}\maketitle

\paragraph{Claim and envelope.}
One fixed cutoff policy consumes a local judgment of the first old delivery
and ordinary permissions at two controlled sites. Its common outputs are an
old fallback and a rescue attempt; a third independently authorized service
continues. The old delivery was owed before the disturbance and is retained
through actual preparation, then either completed at the first cut or
recovered with the preinstalled alternate. A second, localized inability to
execute the first completion may occur while the remote coordinator is
unavailable. The latter contact condition is an authored interface restriction,
not a simulator's discovery of a geographic partition. The plant-side
evidence is complete only in regimes that explicitly supply it. The model
guarantees the named service and uniqueness for all listed source choices
within this final-wave envelope, not for arbitrary mid-wave faults.

\paragraph{Source interpretations.}
The warehouse executes selected RWARE \texttt{step} and grid/cargo code.
A carrier moves the already reserved shelf to the decision junction; a
first scoped motion attempt either reaches a service cell and earns its
native delivery reward or is prevented by the authored unavailable-actuator
cut. The old plant reads both the completed attempt and current occupied
cells. Returning the shelf is still owed even after delivery. The no-
movement alternate uses two explicit turns and an already budgeted move.
The warehouse experiments are measured in action steps; the displayed
$0.75,1,1.25$ indices do not identify them with radio seconds.
The radio fixture executes selected pinned SubT dispatch, radio and callback
functions. Prior sends place old and rescue copies at the alternate sender.
The first attempt may fail its authored RF-path condition; the actual
sender timestamp queue still retains its three-byte charge. The first bounded
invocation has drained at the read cut. Its local tagged application receipt
is used only in the complete-evidence regime. Neither the contact law nor a
site permission is inferred from a task command or display flag.

\paragraph{One cutoff rule and the residual obstruction.}
The fallback, when needed, may use only the right site $R$. The rescue
may use either site, and the unfragmented regime admits one effect per site.
When the first old effect is known to be settled, issue rescue only at an
authorized usable site. If still live and $R$ is usable, reserve $R$ for
the old fallback and use $L$ for rescue if usable. If $R$ is unusable, rescue
alone can still use $L$. When neither status nor the scope of either
conflicting old permission is established, this fixture retains both old
permissions, refuses those two issues and continues the disjoint service.
This last case is not a claim that every masked state needs both sites.
At matched observations and issue laws the native direct enumerator ties.
For both displayed old outcomes with complete evidence, the old and
independent duties settle across both plants. Rescue also settles in the
warehouse; in the radio it can miss if the first attempt failed and the
left debit is live. A direct forced attempt confirms the radio boundary is
a source effect, not just a permission-filter artifact. The source retains
the timestamp from time zero at the exact cutoff $t=1$, expiring it strictly
afterward. The shared decision procedure is byte-identical to its development
freeze. Neither target was independently selected.

\paragraph{Placement of a one-shot gate.}
As a separate adapter extension, suppress a repeat at the actual authority
boundary. An authored warehouse gate at the effect sees the native delivery
queue before committing; if the old shelf is no longer requested it skips
the retry before moving, otherwise it permits the already staged alternate.
It restores all three services without an upstream confirmation report in
this local fixture. The colocated queue is extra evidence compared with an
upstream coordinator that sees only a missing acknowledgement.
In the radio, a terminal once-only callback permits the old retry even when
the sending controller cannot distinguish loss before versus after the first
terminal effect. The retry is authorized by the stipulated original one-shot
contract, not by any newly asserted authority. This restores old and disjoint
service even when acknowledgement is absent. It still charges the alternate
sender three bytes. At early cuts the older receiver's own budget may prevent
the duplicate even reaching the application; at the late cut the application
observes two arrivals and admits one effect. No argument treats application
deduplication as restoring sender airtime.

\paragraph{Alternate granularity and its exact limit.}
With permitted fragmentation as a separate architecture mode, the same
old retry can leave two bytes on each radio sender while a new three-byte
rescue cannot send whole from either. Two ordinary simulator datagrams,
two payload bytes from $L$ and one from $R$, then actually reassemble at
the rescue destination without changing either rolling capacity. Positions
are pre-installed and recognized by origin in this narrow fixture. Extra
header bytes are charged when present. For a payload of $b$ integer bytes,
residual sender budgets $l,r$, and $h$ charged bytes of overhead per nonempty
fragment, there is a feasible two-fragment cut iff
\[
b\ge 2,\qquad l>h,\qquad r>h,\qquad b+2h\le l+r .
\]
Necessity follows by adding the constraints; sufficiency uses the left
payload $\min(l-h,b-1)$ and the positive remainder on the right. At
$t\le 1$ one byte of overhead per fragment defeats that native execution.
The necessary destination-side capacity and all-or-nothing application
reassembly are explicit; sender arithmetic alone does not certify delivery.
This resource can be subdivided; subdivision is not a valid interpretation
of an exclusive warehouse move. The unfragmented and fragmented modes have
different command granularity. The split is a material gain over whole-packet
dispatch at fixed counted payload credits in the stated zero-overhead mode,
not an improvement over an ordinary scheduler allowed to fragment with the
same resources and information.

\paragraph{Common-history calibration and the baseline.}
The source reading route for Contextual Composition is
\url{https://github.com/eternal-november/eos-research-programmes/blob/ed049b09bdaa230afb72bca00d56ad76908c9548/projects/contextual-composition/foundations/native-whole-scope-feedback/Nonblocking_Recurrent_Assembly_Reading_Route.tex}.
A finite diagnostic instance uses the retained-reserve C1 equation of
\texttt{Nonblocking\_History\_Controllers\_and\_Certified\_Exports.tex}.
With equal diagnostic weights on the two native outcomes, the truthful
full-label selection distinguishes fulfilled from still live. At index
one the warehouse's incoming left availability is fair before selection,
then is a point mass (and different) on each selected-label fibre.
Using an unchanged fair reserve after each selection fails C1. At the
same radio index this availability is false in both source histories;
conditioning need not change it. The actual native rows are used, not
independently chosen successful outputs. This is a finite use of the
incoming-reserve test; it does not implement the source's general
probabilistic controller or certify a full imported whole-law gain.

Published NeBula uses local buffering, contact-aware behavior and
rate-limited traffic (\url{https://arxiv.org/html/2504.13461v1},
sections 7 and 8), all relevant comparators. Selected SubT simulator
code is not the NeBula implementation. Application duplicate suppression
and function placement are longstanding, e.g. Saltzer--Reed--Clark,
\url{https://doi.org/10.1145/357401.357402}. Direct exact comparisons tie at
the same permissions for the common core. The known fragmentation
alternative with the same capabilities also ties. The IEEE abstract and
authors' public poster describe HeRo's monitored fault recovery and
PDDL-based workflow, but the full paper is not verified here; no claim is
made about its actual partition, pending authorization or settlement laws.
This is an executable integrated development slice and a precise
substitution boundary, not I1 advancement or an independent holdout.
\end{document}
